
\documentclass[11pt,a4paper]{extarticle}
\usepackage{fontspec}
\usepackage[greek,english]{babel}
\babelfont{rm}{Noto Serif}
\babelfont{sf}{Noto Sans}
\setsansfont{Noto Sans}
\setmainfont{Noto Serif}
\usepackage[a4paper,margin=1.8cm,top=1.65cm,bottom=1.65cm]{geometry}
\usepackage{microtype}
\usepackage{amsmath,amssymb,mathtools}
\usepackage{booktabs,tabularx,array,longtable}
\usepackage{enumitem,xcolor}
\usepackage[most]{tcolorbox}
\usepackage{fancyhdr,hyperref,ragged2e,titlesec,setspace}
\hypersetup{colorlinks=true,linkcolor=black,urlcolor=black,pdftitle={Σημειωματάριο ερευνητικής πρότασης — Μουσική θέση},pdfauthor={Γεώργιος Σπυριδάκης}}
\definecolor{navy}{HTML}{18324A}
\definecolor{slate}{HTML}{4B5D67}
\definecolor{light}{HTML}{F3F6F8}
\definecolor{bluebox}{HTML}{F3F6F8}
\setlength{\parindent}{0pt}
\setlength{\parskip}{5pt}
\linespread{1.035}
\emergencystretch=4em
\sloppy
\setlist[itemize]{leftmargin=1.35em,itemsep=1.5pt,topsep=2pt}
\setlist[enumerate]{leftmargin=1.45em,itemsep=1.5pt,topsep=2pt}
\titleformat{\section}{\sffamily\bfseries\Large\color{navy}}{}{0pt}{}
\titleformat{\subsection}{\sffamily\bfseries\normalsize\color{slate}}{}{0pt}{}
\pagestyle{fancy}
\fancyhf{}
\fancyhead[L]{\sffamily\small Σημειωματάριο ερευνητικής πρότασης}
\fancyhead[R]{\sffamily\small v8.71a / M1--M5}
\fancyfoot[C]{\sffamily\small \thepage}
\newtcolorbox{keybox}[1][]{colback=bluebox,colframe=navy!45!black,boxrule=.45pt,arc=1.5mm,left=2.4mm,right=2.4mm,top=1.6mm,bottom=1.6mm,#1}
\begin{document}
\selectlanguage{greek}
\setcounter{tocdepth}{1}
\begin{center}
{\sffamily\bfseries\LARGE Σημειωματάριο ερευνητικής πρότασης}\\[3mm]
{\sffamily\Large Μουσική θέση, υπολογιστική αναπαράσταση και ιστορική παραλλαγή}\\[3mm]
{\large «Η μουσική θέση ως υπολογιστική μονάδα της βυζαντινής μουσικής»}\\[4mm]
{\small Θεωρητικό πλαίσιο, ερευνητικός σχεδιασμός και κριτήρια ελέγχου · Έκδοση v8.71a}
\end{center}

\begin{keybox}
\textbf{Αντικείμενο.} Η μουσική θέση εξετάζεται ως μεταβλητού μήκους μελωδική μονάδα που μπορεί να αναγνωρίζεται σε διαφορετικά έργα και να αποκτά διαφορετικές πραγματώσεις. Η έρευνα συνδυάζει τη μουσικολογική τεκμηρίωση με τη συμβολική αναπαράσταση, την υπολογιστική αναγνώριση, τη μελέτη ιστορικών εξηγητικών σχέσεων και τον χαρακτηρισμό πλήρων έργων. Η καταλληλότητα της θέσης ως υπολογιστικής μονάδας είναι υπόθεση προς έλεγχο, όχι αρχική βεβαιότητα.
\end{keybox}

\subsection*{Διάρθρωση της έρευνας}
\begin{tabularx}{.97\linewidth}{>{\bfseries\sffamily}p{1.45cm}X}
ΠΕ0 & Πηγές, σώματα δεδομένων, μουσικολογικές υπογραφές οικογενειών, επισημείωση και εφικτότητα.\\
M1 & Επιλογή συμβολικής αναπαράστασης και έλεγχος οριοθέτησης και οικογενειακής αναγνώρισης.\\
M2 & Σύγκριση οικογενειακής με ιστορικά τεκμηριωμένη εξηγητική εποπτεία (B2--EISE).\\
M3a & Ανεξάρτητη μουσικολογική επικύρωση ορίων και οικογενειών.\\
M3c & Επίδραση του πληροφοριακού περιεχομένου των οικογενειών στη μουσική προσδοκία.\\
M3b & Ακροαματική κατηγοριοποίηση συγγενών πραγματώσεων πέρα από την απλή ομοιότητα.\\
M4 & Εκτελεστική μεταβλητότητα και ανεξάρτητο κριτήριο για την ιστορική παραλλαγή.\\
M5 & Χαρακτηρισμός πλήρων έργων ως προς συγκεκριμένο σώμα αναφοράς.\\
\end{tabularx}

\subsection*{Χρήση και καθεστώς του κειμένου}
Το σημειωματάριο συνοψίζει τη μεθοδολογική λογική της πρότασης v8.71a. Δεν αποτελεί πρωτόκολλο με νέα αποτελέσματα. Οι υπό συζήτηση απλοποιήσεις, ιδίως για τη M3b, διακρίνονται από τις επιλογές που ήδη προβλέπει η συγκεκριμένη έκδοση. Πρωτεύουσα πηγή παραμένει το πλήρες κείμενο της ερευνητικής πρότασης.

\clearpage
\tableofcontents
\clearpage


% 1
\section{Το ερευνητικό πρόβλημα και η συνεισφορά}
\subsection*{Η μουσική κλίμακα που ενδιαφέρει την έρευνα}
Η Νέα Μέθοδος αποδίδει τη μελωδία με συνδυασμούς φωνητικών, ποιοτικών, χρονικών και τροπικών σημείων. Η μουσική θέση βρίσκεται σε μια ενδιάμεση κλίμακα ανάμεσα στο μεμονωμένο σημειογραφικό στοιχείο και στο πλήρες έργο. Μπορεί να αναπτύσσεται σε διαφορετικό αριθμό γεγονότων και συλλαβών, χωρίς να χάνει υποχρεωτικά την οικογενειακή της ταυτότητα. Αυτό δημιουργεί πρόβλημα οριοθέτησης, ταξινόμησης και αναπαράστασης της μεταβλητότητας.

\subsection*{Γιατί δεν αρκεί η απλή μελωδική ομοιότητα}
Δύο εμφανίσεις μπορεί να είναι επιφανειακά όμοιες αλλά να επιτελούν διαφορετική τροπική ή φραστική λειτουργία. Αντίστροφα, διαφορετικές αναλυτικές πραγματώσεις μπορεί να ανήκουν στην ίδια οικογένεια. Η μελωδική ομοιότητα, η οικογενειακή συγγένεια και η ιστορική εξηγητική σχέση αντιμετωπίζονται συνεπώς ως τρεις χωριστές κατηγορίες τεκμηρίωσης.

\subsection*{Τι προσθέτει η υπολογιστική προσέγγιση}
Η έρευνα επιδιώκει να αναγνωρίζει σε νέα έργα τις εμφανίσεις θέσεων, να εκτιμά πιθανή οικογενειακή ένταξη και να καταγράφει την αβεβαιότητα. Παράλληλα εξετάζει ποια πληροφορία από ιστορικές εξηγήσεις μπορεί να αξιοποιηθεί κατά την εκμάθηση των μουσικών αναπαραστάσεων. Ο τελικός στόχος είναι η συγκρίσιμη περιγραφή πλήρων έργων με διατήρηση της σύνδεσης προς τις συγκεκριμένες πηγές και τα συγκεκριμένα χωρία.

\subsection*{Η μεθοδολογική αφετηρία}
Προηγούμενες υπολογιστικές εργασίες έχουν εξετάσει τη μουσική τμηματοποίηση, την ανεύρεση επαναλαμβανόμενων σχημάτων και τη μελωδική ομοιότητα. Η εργασία HarmonyToC αποτελεί χρήσιμο προηγούμενο για την πειραματική επιλογή της συμβολικής μονάδας. Δεν αποτελεί όμως έτοιμο μοντέλο οικογενειών βυζαντινών θέσεων. Στη M1 η επιλογή αναπαράστασης ελέγχεται με κοινές δοκιμασίες αντί να καθορίζεται μόνο από τη διαισθητική καταλληλότητά της.

\begin{keybox}
\textbf{Κεντρική διάκριση.} Η ιστορική μουσικολογία ορίζει τι χρειάζεται να εξεταστεί. Η υπολογιστική μέθοδος προτείνει μετρήσιμο τρόπο αναπαράστασης και αναγνώρισης. Η ανεξάρτητη αξιολόγηση ελέγχει κατά πόσο η προτεινόμενη λύση ανταποκρίνεται στο μουσικό αντικείμενο.
\end{keybox}


% PAGE 3
\section{Μουσικολογική βάση: πέντε έννοιες, τρεις σχέσεις}
\subsection*{Η βασική οντολογία}
\begin{tabularx}{\linewidth}{>{\bfseries}p{3cm}X}
Θέση (τύπος) & Μουσικός σχηματισμός με διαγνωστικές σχέσεις, μεταβλητή έκταση και τροπική/φραστική λειτουργία.\tabularnewline
Οικογένεια & Επιχειρησιακή αναλυτική κατηγορία συγγενών πραγματώσεων· δεν είναι αυτόματα ιστορική ονομασία.\tabularnewline
Εμφάνιση & Συγκεκριμένο τμήμα ενός έργου με όρια, προέλευση και ενδεχόμενη αβεβαιότητα.\tabularnewline
Πραγμάτωση & Η συγκεκριμένη σημειογραφική/μελωδική μορφή που έχει μια θέση σε αυτή την εμφάνιση.\tabularnewline
Ιστορική εξήγηση & Σχέση συγκεκριμένων πηγών/αποδόσεων με αποδείξιμη ιστορική προέλευση· όχι απλή ομοιότητα.\tabularnewline
\end{tabularx}

\subsection*{Τρία διαφορετικά είδη σχέσεων}
\[E_{sim}\neq E_{fam}\neq E_{ex}.\]
Η μελωδική ομοιότητα δεν αρκεί για οικογενειακή ένταξη, και η ίδια οικογένεια δεν σημαίνει ιστορικά τεκμηριωμένη εξηγητική σχέση. Οι ακμές εξήγησης απαιτούν πηγή, κοινό χωρίο και έλεγχο τεκμηρίων.

\subsection*{Ιστορική θεμελίωση: τι αποδίδουμε με ασφάλεια}
Ο Μανουήλ Χρυσάφης και ο Γαβριήλ ιερομόναχος αποτελούν ιστορικά σημεία αναφοράς για τη θέση ως οργανωμένο μελικό σχηματισμό, όχι για τη συγκεκριμένη σημερινή υπολογιστική οικογένεια. Στον Χρύσανθο, ειδικά στο \emph{Θεωρητικόν Μέγα} §408, η εξάρτηση της ενέργειας παλαιότερων σχηματισμών από τη βαθμίδα εφαρμογής ενισχύει τη σημασία του συμφραζομένου· δεν αποτελεί προδιαγραφή του EISE ή του PPM. Η κριτική εξέταση των πηγών είναι αναγκαία όταν ισχυριζόμαστε εξηγητική σχέση.

\subsection*{Γιατί μεταβλητό μήκος;}
Η θέση δεν συμπίπτει υποχρεωτικά με ένα νεύμα, ένα ατομικό γεγονός, μία συλλαβή ή ολόκληρη φράση. Μπορεί να υπάρχουν συνδετικές και ανεπίλυτες περιοχές, καθώς και τεκμηριωμένες επικαλύψεις. Η βασική έξοδος δεν επιβάλλει ετικέτα οικογένειας σε κάθε γεγονός.




% PAGE 4
\section{ΠΕ0: πηγές, κάρτες οικογενειών, επισημείωση}
\subsection*{Τα δεδομένα έχουν διαφορετικούς ρόλους}
\begin{itemize}
\item \textbf{Κατάλογος}: η παπαδική συλλογή της Μονής Αγίου Αντωνίου καλύπτει υλικό προέλευσης από 758 συνθέσεις· είναι καταλογογραφική αφετηρία, όχι αυτόματη βάση πλήρως επισημειωμένων θέσεων. Οι 10.000+ φόρμουλες ανήκουν σε διαφορετική, ειρμολογική/στιχηραρική συλλογή.
\item \textbf{Μητρώο εμφανίσεων}: έργο, πηγή, ήχος, κειμενικό χωρίο, όρια, λειτουργία, ανεπίλυτο υλικό, επιτρεπτές εναλλακτικές, προέλευση.
\item \textbf{Ιστορικό μητρώο σχέσεων}: χωριστή σχέση εξήγησης, κοινό χωρίο, ανεξαρτησία και εξάρτηση πηγών. Δεν παράγουμε θετικά ζεύγη ιστορικής εξήγησης συνθετικά.
\item \textbf{Ανεξάρτητη αξιολόγηση}: έργα και πηγές εκτός εκπαίδευσης/επιλογής, ώστε να μην υπάρχει διαρροή.
\end{itemize}
\subsection*{Κάρτα οικογένειας: τι ορίζουμε πριν μάθει το μοντέλο}
Αναλυτική κλίμακα, αναγκαίες μελωδικές σχέσεις, επιτρεπτές μεταβολές, τροπικοί περιορισμοί όπου τεκμηριώνονται, θετικά και δύσκολα αρνητικά παραδείγματα, κανόνες μη ένταξης και δυνατότητα \emph{ανεπίλυτο}. Οι πρώτες υποψήφιες περιπτώσεις (π.χ. παρακλητική, κύλισμα, πίεσμα) δεν θεωρούνται αυτόματα επικυρωμένες οικογένειες.
\subsection*{Διακριτές φάσεις επισημείωσης}
\begin{tabularx}{\linewidth}{>{\bfseries}p{3cm}X}
ΠΕ0 / ανάπτυξη & Εσύ και δεύτερος ανεξάρτητος ειδικός σε προκαθορισμένο μέρος· δημιουργία και έλεγχος εγχειριδίου και δεδομένων ανάπτυξης.\tabularnewline
Πιλότος οδηγιών & 2--4 ενεργοί ψάλτες, χωριστοί από τη M3a· ελέγχουν εάν η εργασία είναι κατανοητή. Οι ετικέτες τους \emph{δεν} γίνονται επιβεβαιωτικά δεδομένα.\tabularnewline
M3a / τυφλή μελέτη & 3--5 άλλοι ανεξάρτητοι ειδικοί, σε χωριστό σύνολο, χωρίς προβλέψεις μοντέλων και χωρίς τις κρίσεις των άλλων.\tabularnewline
\end{tabularx}



% 4
\section{M1 — Συμβολικές μονάδες και επίπεδα πληροφορίας}
\subsection*{Η βασική ερώτηση}
Η M1 εξετάζει με ποια κλίμακα συμβολικής αναπαράστασης μπορούν να εντοπιστούν αποτελεσματικότερα τα όρια μιας θέσης και να αναγνωριστεί η οικογενειακή της ταυτότητα. Δεν προϋποθέτει ότι η μεγαλύτερη ή η πιο σύνθετη μονάδα είναι πάντοτε καταλληλότερη.

\subsection*{Οι τρεις κλίμακες της μονάδας εισόδου}
\begin{tabularx}{\linewidth}{>{\bfseries}p{1.25cm}p{4cm}X}
\(u_1\) & Σημειογραφικός χαρακτήρας & Κλίμακα πλησιέστερη στη γραπτή μορφή της πηγής.\\
\(u_2\) & Ατομικό μελωδικό γεγονός & Κοινό λεπτομερές επίπεδο αναφοράς· δεν ισοδυναμεί πάντοτε με ένα νεύμα.\\
\(u_3\) & Σύνθετη νευματική μονάδα & Προκαθορισμένη ομαδοποίηση μελωδικών γεγονότων με σταθερό κανόνα.\\
\end{tabularx}
Η ομαδοποίηση \(u_3=f_{\mathrm{group}}(N)\) παραμένει ίδια σε όλες τις δοκιμασίες, ώστε να μη συγχέεται η αλλαγή μονάδας με την προσθήκη νέων γνωρισμάτων.

\subsection*{Τα τέσσερα σωρευτικά επίπεδα πληροφορίας}
\[
F=\{N,\ N+T,\ N+T+M,\ N+T+M+L\}.
\]
Το \(N\) περιλαμβάνει σημειογραφική και μελωδική πληροφορία, το \(T\) χρονική, το \(M\) τροπική και το \(L\) κειμενική και προσωδιακή. Ο διασταυρωμένος σχεδιασμός έχει τρεις κλίμακες μονάδας και τέσσερις προκαθορισμένες βαθμίδες γνωρισμάτων: συνολικά δώδεκα κύριες συνθήκες.

\subsection*{Τι ακριβώς εξετάζει η προσθήκη πληροφοριών}
Οι μεταβάσεις από \(N\) σε \(N+T\), έπειτα σε \(N+T+M\) και τέλος σε \(N+T+M+L\) εκτιμούν τη συνεισφορά κάθε νέου συνόλου \emph{υπό τα προηγούμενα ήδη διαθέσιμα γνωρίσματα}. Δεν επιτρέπουν αυτομάτως ανεξάρτητο συμπέρασμα για τη συνολική επίδραση του χρόνου, της τροπικότητας ή του κειμένου. Πρόσθετες επιμέρους αφαιρέσεις γνωρισμάτων προορίζονται για ερμηνεία και έλεγχο ανθεκτικότητας, όχι για νέα επιλογή της κύριας αναπαράστασης.


% 5
\section{M1 — Δίκαιη σύγκριση και επιλογή της \(R^\star\)}
\subsection*{Κοινός στόχος για διαφορετικές αναπαραστάσεις}
Οι τρεις κλίμακες της M1 δεν έχουν το ίδιο μήκος ακολουθίας ούτε το ίδιο λεξιλόγιο συμβόλων. Για συγκρίσιμη αξιολόγηση, όμως, η πρόβλεψη αφορά το ίδιο επόμενο ατομικό μελωδικό γεγονός και τα όρια εξετάζονται στην ίδια κοινή κλίμακα γεγονότων. Οι διαφορές εισόδου δεν επιτρέπεται να αλλάζουν ανεξέλεγκτα το μουσικό αντικείμενο της δοκιμασίας.

\subsection*{Δύο κριτήρια με διαφορετικό ρόλο}
Το \(M(R)\) αποτυπώνει την επίδοση εντοπισμού ορίων με τον δείκτη F1 και την προκαθορισμένη ανοχή. Το \(G(R)\) αποτυπώνει τη μακρο-μέση ανάκληση των οικογενειών: κάθε οικογένεια έχει συγκρίσιμο βάρος ανεξάρτητα από το πλήθος των εμφανίσεών της.

\subsection*{Ελάχιστο όριο οικογενειακής επάρκειας}
\[
g_{\min}=\max\{g_{\mathrm{use}},\widehat G_0^{\mathrm{PE0}}+\delta_{\mathrm{use}}\}.
\]
Το όριο ορίζεται πριν εξεταστούν οι κύριες συνθήκες. Οι επιλέξιμες αναπαραστάσεις πρέπει πρώτα να καλύπτουν την ελάχιστη απαίτηση οικογενειακής αναγνώρισης. Μεταξύ αυτών επιλέγεται εκείνη με την καλύτερη επίδοση οριοθέτησης:
\[
\widehat{\mathcal A}_{\mathrm{dec}}=\{R:\widehat G(R)\ge g_{\min}\},\qquad
R^\star=\arg\max_{R\in\widehat{\mathcal A}_{\mathrm{dec}}}\widehat M(R).
\]

\subsection*{Η απόφαση δεν σημαίνει βεβαιότητα υπεροχής}
Η \(R^\star\) είναι η λειτουργική επιλογή με την οποία συνεχίζεται η M2. Δεν σημαίνει ότι έχει αποδειχθεί πως μία αναπαράσταση είναι μοναδικά καλύτερη σε κάθε έργο. Αναφέρονται επίσης η αβεβαιότητα επιλογής \(\mathcal R^\star\) και η ενδεχόμενη απώλεια στην ανάκληση οικογενειών:
\[
\Delta_G^{\mathrm{loss}}=G(R^\star)-\max_R G(R)\le 0.
\]
Εάν καμία συνθήκη δεν πληροί το ελάχιστο όριο, δεν μειώνεται εκ των υστέρων η απαίτηση για να καταστεί δυνατή η κύρια επιβεβαιωτική σύγκριση.


% 6
\section{M2 — Το πρόβλημα αναγνώρισης τμημάτων}
\subsection*{Η έξοδος που απαιτεί το μουσικό αντικείμενο}
Η είσοδος είναι μία συμβολική ακολουθία μελωδικών γεγονότων στην αναπαράσταση \(R^\star\). Η επιδιωκόμενη έξοδος είναι σύνολο εμφανίσεων, καθεμία με αρχή, τέλος, πιθανή οικογένεια και εκτίμηση αβεβαιότητας. Μπορεί να υπάρχουν συνδετικές ή ανεπίλυτες περιοχές, χωρίς υποχρεωτική οικογενειακή ετικέτα σε κάθε γεγονός.

\subsection*{Γιατί χρειάζεται ειδικό μοντέλο τμημάτων}
Οι μουσικές θέσεις έχουν μεταβλητή έκταση. Επομένως, δεν αρκεί ένας ταξινομητής που αποδίδει ανεξάρτητη ετικέτα σε κάθε μεμονωμένο γεγονός. Το ημιμαρκοβιανό υπόδειγμα CRF βαθμολογεί \emph{ολόκληρα υποψήφια τμήματα} και επιλέγει από κοινού τα όριά τους και την πιθανή οικογένεια. Η αρχική μορφή μηδενικής τάξης δεν βασίζεται σε σχέσεις διαδοχής μεταξύ οικογενειών για να κερδίζει τεχνητά την αναγνώρισή τους.

\subsection*{Μερικές επισημειώσεις και αποδεκτές εναλλακτικές}
Σε πραγματικές πηγές δεν είναι γνωστές πάντοτε όλες οι θέσεις. Μία περιοχή χωρίς επισημείωση δεν αποτελεί αναγκαστικά αρνητικό παράδειγμα. Η συνάρτηση εκπαίδευσης επιτρέπει ένα σύνολο συμβατών αναλύσεων, αντί να απαιτεί μία μοναδική τελική ετικέτα εκεί όπου τα τεκμήρια παραμένουν αβέβαια.

\subsection*{Άγνωστες οικογένειες και αποχή}
Μία αναγνωρισμένη εμφάνιση μπορεί να μη συγκρίνεται με ασφάλεια με το γνωστό λεξιλόγιο. Το σύστημα προβλέπει δυνατότητα αποχής ή καταγραφής αγνώστου οικογένειας. Αυτό δεν ισοδυναμεί με απόδειξη ότι ανακαλύφθηκε ιστορικά νέος τύπος. Απαιτείται ανεξάρτητη μουσικολογική διερεύνηση.

\subsection*{Η αρχιτεκτονική ως ελέγξιμη επιλογή}
Η ημιμαρκοβιανή διαδικασία αποτελεί το βασικό σχεδιαζόμενο υπόδειγμα, αλλά συγκρίνεται στο στάδιο ανάπτυξης και με απλούστερες προσεγγίσεις. Εάν η πρόσθετη πολυπλοκότητα δεν προσφέρει μετρήσιμη αξία, μπορεί να προτιμηθεί απλούστερη λύση με διατήρηση του ίδιου ερευνητικού ερωτήματος.


% 7
\section{EISE — Αρχιτεκτονική και αντικειμενική συνάρτηση}
\subsection*{Τα βασικά στάδια επεξεργασίας}
Ο κωδικοποιητής μετατρέπει κάθε ατομικό μελωδικό γεγονός σε αριθμητική κατάσταση \(h_t\). Από τις καταστάσεις αυτές σχηματίζονται υποψήφια συνεχόμενα τμήματα. Για κάθε τμήμα δημιουργείται συνοπτική αναπαράσταση \(r_s\) και, έπειτα, \(z_{\mathrm{id}}\), που χρησιμοποιείται για την οικογενειακή ταυτότητα. Οι τοπικές καταστάσεις \(h_t\) διατηρούν πληροφορία για τις ιδιαίτερες μορφές πραγμάτωσης.
\[
r_s=\left[\frac{1}{b-a+1}\sum_{t=a}^{b}h_t;\ h_a;\ h_b\right],\qquad
z_{\mathrm{id}}(s)=\operatorname{Normalize}(q_{\mathrm{id}}(r_s)).
\]
Η αντιστοίχιση με μια γνωστή οικογένεια εκτιμάται, μεταξύ άλλων, από την απόσταση της αναπαράστασης του τμήματος από ένα πρότυπο οικογένειας. Η τελική διαδικασία επιλέγει ταυτόχρονα όρια και οικογενειακές αποδόσεις.

\subsection*{Η συνάρτηση εκπαίδευσης}
\[
\mathcal L_{\mathrm{EISE}}=
\mathcal L_{X,\mathrm{seg}}^{\mathrm{partial}}
+\lambda_{\mathrm{rel}}\mathcal L_{\mathrm{pairs}}^{\mathrm{id}}
+\lambda_{\mathrm{proto}}\mathcal L_{\mathrm{proto}}
+\lambda_{\mathrm{loc}}\mathcal L_{\mathrm{local}}.
\]
Ο πρώτος όρος αφορά τα όρια και τις οικογενειακές ετικέτες με δυνατότητα μερικής επισημείωσης. Ο δεύτερος οργανώνει την ταυτότητα με βάση επιλεγμένα θετικά ζεύγη. Ο τρίτος ενισχύει τη συνοχή των γνωστών οικογενειών γύρω από χαρακτηριστικές αναπαραστάσεις. Ο τέταρτος αποτρέπει την απώλεια πληροφορίας για τη συγκεκριμένη μελωδική πραγμάτωση.

\subsection*{Γιατί διατηρούνται δύο επίπεδα αναπαράστασης}
Η οικογενειακή αναπαράσταση πρέπει να αναγνωρίζει το κοινό στοιχείο ανάμεσα σε διαφορετικές πραγματώσεις, χωρίς να συγχέει τις οικογένειες. Ταυτόχρονα, η μελέτη της ιστορικής και εκτελεστικής μεταβλητότητας απαιτεί πληροφορία για τις τοπικές διαφοροποιήσεις. Γι' αυτό το \(z_{\mathrm{id}}\) και οι \(h_t\) έχουν διαφορετικούς, συμπληρωματικούς ρόλους.

\subsection*{Τι παραμένει προς έλεγχο}
Το συγκεκριμένο πλήθος επιπέδων, οι διαστάσεις των αναπαραστάσεων, οι υπερπαράμετροι και η αξιοπιστία της διατήρησης τοπικής πληροφορίας θα ελεγχθούν στην πιλοτική φάση. Η αρχιτεκτονική αποτελεί ερευνητική πρόταση, όχι ήδη επιβεβαιωμένο αποτέλεσμα.


% 8
\section{M2 — Δίκαιη σύγκριση B0, B1, B2 και EISE}
\begin{tabularx}{\linewidth}{>{\bfseries}p{1.35cm}X}
B0 & Άμεσος συμβολικός χαρακτηρισμός της ιστορικής παραλλαγής, χωρίς εκμάθηση αριθμητικής αναπαράστασης.\\
B1 & Κοινή νευρωνική και τμηματική υποδομή, χωρίς την πρόσθετη εποπτεία από ζεύγη πραγματώσεων.\\
B2 & Ίδια υποδομή με εποπτεία από ζεύγη της ίδιας οικογένειας χωρίς τεκμηριωμένη ιστορική εξηγητική σχέση.\\
EISE & Ίδια υποδομή, αλλά η σχεσιακή εποπτεία προέρχεται από ανεξάρτητα τεκμηριωμένες ιστορικές εξηγητικές σχέσεις.\\
\end{tabularx}

\subsection*{Η κεντρική αντίθεση}
\[
\Delta_{\mathrm{hist}}=S_{\mathrm{EISE}}-S_{B2},\qquad
\Delta_{\mathrm{learn}}=S_{\mathrm{EISE}}-S_{B0}.
\]
Το \(S_m\) αποτυπώνει κατά πόσο το ιστορικό προφίλ παραλλαγής του αντίστοιχου μοντέλου συμφωνεί με ένα ανεξάρτητο σύγχρονο εκτελεστικό προφίλ. Η κύρια σύγκριση \(\Delta_{\mathrm{hist}}\) εξετάζει την πρόσθετη αξία της \emph{προέλευσης της σχέσης εποπτείας}, όχι την ιστορική εξήγηση ως αιτία μιας σημερινής εκτέλεσης.

\subsection*{Προϋποθέσεις δίκαιης αντιπαραβολής}
Οι B2 και EISE πρέπει να έχουν κοινή αρχιτεκτονική, κοινές οικογενειακές ετικέτες, συγκρίσιμο πλήθος θετικών ζευγών, αντίστοιχη δυσκολία παραδειγμάτων, συγκρίσιμη έκθεση σε δεδομένα και ίδιους διαθέσιμους υπολογιστικούς πόρους. Η διαφοροποίηση που επιδιώκεται να απομονωθεί είναι η τεκμηριωμένη ιστορική προέλευση των ζευγών της EISE.

\subsection*{Ερμηνεία πιθανών αποτελεσμάτων}
Μια θετική \(\Delta_{\mathrm{hist}}\), με την απαιτούμενη στατιστική υποστήριξη, θα ενίσχυε την αξία της ιστορικής εποπτείας υπό το συγκεκριμένο εξωτερικό κριτήριο. Τιμή κοντά στο μηδέν δεν θα τεκμηρίωνε πρόσθετη βελτίωση, χωρίς να ακυρώνει τη μουσικολογική σημασία της εξήγησης. Αρνητική τιμή θα απαιτούσε εξέταση της κάλυψης, της ετερογένειας των πηγών και των ορίων του μοντέλου. Η B2 δεν θεωρείται εκ των προτέρων ανίκανη να μάθει εσωτερική σταθερότητα πραγματώσεων.


% PAGE 10
\section{M3a — ανεξάρτητη μουσικολογική επικύρωση}
\subsection*{Διαφορετικός στόχος από το ΠΕ0}
Το ΠΕ0 κατασκευάζει και ελέγχει το λεξιλόγιο ανάπτυξης. Η M3a ελέγχει αν άλλοι γνώστες της παράδοσης εφαρμόζουν το ίδιο μουσικολογικό πλαίσιο σε δεδομένα που δεν χρησιμοποιήθηκαν στην ανάπτυξη των μοντέλων. Η M3a μπορεί να προχωρά παράλληλα με την M2, διατηρώντας ανεξάρτητα δεδομένα.
\subsection*{Η τυφλή μελέτη}
\begin{itemize}
\item 3--5 ανεξάρτητοι ειδικοί με επιδίωξη ποικιλίας ψαλτικής/ερευνητικής προέλευσης.
\item Βλέπουν το αναγκαίο μουσικό συμφραζόμενο και τις οριστικοποιημένες οδηγίες, \emph{όχι} τις προβλέψεις του EISE ή τις κρίσεις των άλλων αξιολογητών.
\item Επισημειώνουν όρια, οικογενειακή ταυτότητα και, όπου τεκμηριώνεται, λειτουργικό ρόλο· έχουν επιλογή \emph{ανεπίλυτο}.
\item Μετράμε χωριστά συμφωνία για όρια και για οικογενειακή απόδοση, μαζί με την έκταση των ανεπίλυτων περιπτώσεων.
\item Δεν αλλάζουμε αναδρομικά την προδιαγραφή του μοντέλου ή το σύνολο τελικής αξιολόγησης ώστε να αυξηθεί η συμφωνία.
\end{itemize}
\subsection*{Τι σημαίνει καλό αποτέλεσμα;}
Ο συνδυασμός συμφωνίας μεταξύ ειδικών, ανεξάρτητης κριτικής των πηγών και καλής ανάκτησης σε νέα έργα στηρίζει τη χρησιμότητα του λεξιλογίου. Δεν τεκμηριώνει ότι κάθε οικογένεια φέρει μία ενιαία ιστορική ονομασία. Η διαφωνία μπορεί να δείχνει πραγματική μεταβλητότητα της παράδοσης.



% PAGE 11
\section{M3c — μουσική προσδοκία, PPM και ανθρώπινη κρίση}
\subsection*{Ερευνητικό ερώτημα}
Ο ακροατής ακούει ένα μουσικό απόσπασμα και τη συνέχεια και βαθμολογεί πόσο απρόσμενη του φαίνεται. Το ζητούμενο είναι αν η οικογένεια της συνέχειας έχει πρόσθετη ερμηνευτική αξία σε σχέση με τα μεμονωμένα γεγονότα.
\subsection*{Δύο διαφορετικά πληροφοριακά μεγέθη}
\[
I_F=-\log_2 P(F_t\mid F_{<t},M,L),\qquad I_{event}=\text{σύνοψη της έκπληξης των γεγονότων}.
\]
Το \(I_F\) αφορά πόσο απρόσμενη είναι η οικογένεια \emph{μετά την προηγούμενη διαδοχή} και με δεδομένο τροπικό/δομικό/κειμενικό περιβάλλον. Το \(I_{\mathrm{event}}\) αφορά τα μεμονωμένα γεγονότα της ίδιας συνέχειας, με κύρια σύνοψη το μέσο πληροφοριακό περιεχόμενό τους. Δεν συγκρίνουμε απλώς τις ακατέργαστες αριθμητικές τιμές τους.
\subsection*{Γιατί PPM και ποιος ο υπολογιστικός έλεγχος;}
Το PPM (πρόβλεψη από μερική αντιστοίχιση) χρησιμοποιεί προηγούμενο συμφραζόμενο μεταβλητού μήκους και, όταν αυτό είναι σπάνιο, επιστρέφει σε μικρότερο. Ελέγχεται σε έργα εκτός εκπαίδευσης έναντι του μοντέλου \(P(F_t\mid M,L)\), ώστε να φανεί αν \emph{η ακολουθία} προσφέρει πληροφορία πέρα από τον ήχο/δομικό τόπο.
\subsection*{Γιατί υψηλό/χαμηλό στους αντίθετους συνδυασμούς;}
Χρειαζόμαστε περιπτώσεις \(I_F\uparrow,I_{event}\downarrow\) και \(I_F\downarrow,I_{event}\uparrow\), όχι για να «πάμε στα άκρα», αλλά για να χωρίσουμε τις επιδράσεις. Η μουσική φυσικότητα ελέγχεται τυφλά. Αν οι περιπτώσεις δεν επαρκούν, η κύρια υπόθεση δεν εξαναγκάζεται.
\subsection*{Εμπειρία και λατρευτικό πλαίσιο}
Ο \(E_{Byz}\) είναι διαβαθμισμένος δείκτης ψαλτικής εξοικείωσης, όχι απλώς ψάλτης/μη ψάλτης. Το \(R\) είναι πειραματική \emph{μη μουσική ένδειξη} λατρευτικού πλαισίου, όχι αναπαράσταση ολόκληρης της επιτέλεσης. Το \(\Delta_{ctx}\) εξετάζει αν αλλάζει διαφορετικά η σχέση των κρίσεων με \(I_F\) και \(I_{\mathrm{event}}\).


% PAGE 12
\section{M3b — κατηγοριοποίηση πέρα από επιφανειακή ομοιότητα}
\subsection*{Τρία αποσπάσματα, ένα ερώτημα}
Ο ακροατής ακούει τρεις πραγματώσεις και επιλέγει ποια ταιριάζει λιγότερο με τις άλλες δύο. Σε μεικτή τριάδα έχουμε δύο πραγματώσεις της ίδιας οικογένειας και μία άλλης. Περιλαμβάνονται επίσης τριάδες με τρεις πραγματώσεις της ίδιας οικογένειας, για τη μελέτη της εσωτερικής διαφοροποίησης.
\subsection*{Το ερευνητικό πρόβλημα}
Αν η οικογενειακή συγγένεια συμπίπτει πάντα με τη μεγαλύτερη ακουστική ομοιότητα, μια σωστή απάντηση δεν αποδεικνύει κατηγοριοποίηση. Γι' αυτό επιλέγουμε τριάδες από \emph{κοινή περιοχή συμβολικών και ακουστικών αποστάσεων} και περιλαμβάνουμε δύσκολα παραδείγματα: ίδια οικογένεια με μεγάλη επιφανειακή διαφορά, διαφορετικές οικογένειες με συγκρίσιμη επιφανειακή ομοιότητα.
\subsection*{Τι εξετάζει η υπόθεση \(H_{\mathrm{cat}}\);}
Εξετάζουμε αν η οικογενειακή ένταξη προβλέπει την ανθρώπινη επιλογή \emph{αφού ελεγχθούν} οι επιφανειακές αποστάσεις και άλλες προσυμφωνημένες μεταβλητές. Μια θετική σχέση δεν αποδεικνύει ότι ο ακροατής γνωρίζει ιστορικές ονομασίες ή εξηγητικές σχέσεις.
\subsection*{Καθεστώς στην έκδοση v8.71a}
\textbf{Στη v8.71a η M3b προβλέπεται ως επιβεβαιωτική δοκιμασία} με προκαθορισμένα κριτήρια επάρκειας, και ακολουθεί τη M3c στην ανθρώπινη συλλογή. Η πιθανή μεταφορά της σε προαιρετική βαθμίδα είναι ενδεχόμενο μελλοντικής αναθεώρησης, όχι ήδη εγκεκριμένη αλλαγή. Αν γίνει, αλλάζουν αντίστοιχα προκαταγραφή του ερευνητικού πρωτοκόλλου, στατιστική ισχύς και οικογένεια ελέγχων.



% PAGE 13
\section{M4 — γραφή, εκτέλεση και ανεξάρτητο κριτήριο}
\subsection*{Πρώτα αντιστοίχιση, έπειτα συμπέρασμα}
Χρειαζόμαστε \textbf{πολλαπλές ηχογραφήσεις του ίδιου συγκρίσιμου μουσικού υλικού}. Ευθυγραμμίζουμε χειροκίνητα τα όρια συλλαβών και, όπου η ακρίβεια το επιτρέπει, αντιστοιχίζουμε εντός της συλλαβής τα επιμέρους μελωδικά γεγονότα με ευθυγράμμιση δυναμικής στρέβλωσης χρόνου με καθοδήγηση από την παρτιτούρα ή μοντέλο καταστάσεων. Μια μελισματική συλλαβή δεν είναι ένα γεγονός. Αν η λεπτομερής ευθυγράμμιση αποτύχει, εργαζόμαστε σε συλλαβή/νευματική ομάδα.
\subsection*{Διάκριση χρονισμού και μελωδικής ανάπτυξης}
Μια πρόσθετη μελωδική ανάπτυξη δεν πρέπει να απορροφηθεί ως «καθυστέρηση» σε έναν απλό χρονικό μετασχηματισμό. Την επισημαίνουμε ως δομική/μελωδική διαφοροποίηση όταν τα διαθέσιμα δεδομένα την τεκμηριώνουν. Χωριστά καταγράφουμε κανονικοποιημένο χρονισμό, συνολική αγωγή, τοπική αγωγή, καταληκτική επιμήκυνση και, όπου αξιόπιστο, τονικό ύψος.
\subsection*{Η κύρια εξωτερική σύγκριση της M2}
\begin{itemize}
\item \(A(z)\): ιστορικό προφίλ παραλλαγής ανά εσωτερική περιοχή \(z\) της θέσης.
\item \(\Gamma(z)\): ανεξάρτητη διασπορά μεταξύ σύγχρονων εκτελεστών στην ίδια μουσική περιοχή.
\item \(S_m\): συμφωνία κατάταξης ανάμεσα σε ιστορικό και εκτελεστικό προφίλ για \(m\in\{B0,B2,EISE\}\).
\end{itemize}
Η προεπιλεγμένη κύρια εκτελεστική μεταβλητή είναι η διασπορά \emph{κανονικοποιημένου λογαριθμικού χρονισμού}, αν η ποιότητά της επαρκεί. Το ύψος είναι δευτερεύον. Τα εκτελεστικά δεδομένα δεν αναπροσαρμόζουν εκ των υστέρων τον EISE.
\subsection*{Όριο ιστορικού ισχυρισμού}
Η σύμπτωση μιας σύγχρονης εκτελεστικής μορφής με μία ιστορική εξήγηση είναι μορφολογική/εκτελεστική αντιστοιχία· \emph{δεν} αποδεικνύει συγκεκριμένη διδασκαλική γενεαλογία ή ανακατασκευή του ιστορικού ακούσματος.


% 13
\section{M5 — Από τις εμφανίσεις στο πλήρες μουσικό έργο}
\subsection*{Η θέση μέσα στην ιεραρχία της σύνθεσης}
Η θέση δεν αντικαθιστά τα επιμέρους μελωδικά γεγονότα ούτε το επίπεδο ολόκληρου του έργου. Συνδέει τις δύο κλίμακες. Οι αναγνωρισμένες εμφανίσεις συγκροτούν χάρτη θέσεων στον οποίο διατηρούνται τα όρια, οι οικογένειες, η εσωτερική πραγμάτωση, τα τυχόν συνδετικά χωρία και η προέλευση από τη μουσική πηγή.

\subsection*{Τέσσερις ανεξάρτητες διαστάσεις χαρακτηρισμού}
\[
\mathbf D(w;\Omega)=\bigl(D_{\mathrm{fam}},D_{\mathrm{seq}},D_{\mathrm{func}},D_{\mathrm{within}}\bigr).
\]
\begin{itemize}
\item \(D_{\mathrm{fam}}\): διαφορές στη συχνότητα και την κατανομή των οικογενειών.
\item \(D_{\mathrm{seq}}\): διαφορές στη σειρά με την οποία εμφανίζονται οι οικογένειες.
\item \(D_{\mathrm{func}}\): διαφορές στη δομική, τροπική ή κειμενική τοποθέτησή τους.
\item \(D_{\mathrm{within}}\): διαφορές στον τρόπο πραγμάτωσης εμφανίσεων της ίδιας οικογένειας.
\end{itemize}

\subsection*{Γιατί δεν αρκεί μία συνολική βαθμολογία}
Ένα έργο μπορεί να χρησιμοποιεί συνηθισμένες οικογένειες σε ασυνήθιστη σειρά. Άλλο έργο μπορεί να διατηρεί τη συνήθη σειρά αλλά να αναπτύσσει εκτενέστερα τις ίδιες θέσεις. Εάν οι διαφορές αυτές συγχωνευθούν, χάνεται η μουσικολογική τους ερμηνεία. Η M5 παράγει συνεπώς πολύπλευρη περιγραφή, όχι βαθμό ορθότητας, αισθητικής αξίας ή «κλασικότητας».

\subsection*{Αναγωγή από το αποτέλεσμα στην πηγή}
Κάθε αριθμητικό εύρημα πρέπει να συνδέεται με συγκεκριμένο έργο, οικογένεια, εμφάνιση και χωρίο της πηγής. Η ενδοοικογενειακή απόσταση \(D_{\mathrm{within}}\) παραμένει ενιαία και συνοδεύεται από επιμέρους ερμηνεύσιμους δείκτες μόνο όπου το υλικό επιτρέπει αξιόπιστη σύγκριση.


% 14
\section{M5 — Ιστορικό περιβάλλον αναφοράς και όρια ισχυρισμών}
\subsection*{Τι περιλαμβάνει το περιβάλλον αναφοράς}
\[
\Omega=(D_\Omega,\rho,\Pi,\mathcal K_\Omega).
\]
Το \(D_\Omega\) είναι το δηλωμένο σώμα μουσικών πηγών. Το \(\rho\) περιλαμβάνει την οριστικοποιημένη αναπαράσταση, το λεξιλόγιο οικογενειών και τους κανόνες αποχής. Το \(\Pi\) είναι το υπόδειγμα πιθανοτικής διαδοχής οικογενειών, ενώ το \(\mathcal K_\Omega\) συγκεντρώνει τους κανόνες ιστορικής ένταξης, χρονολόγησης και ανεξαρτησίας των πηγών. Επομένως, κάθε διαπίστωση διαφοροποίησης εξαρτάται από το συγκεκριμένο περιβάλλον αναφοράς.

\subsection*{Ελάχιστη και εκτεταμένη ιστορική εφαρμογή}
Η προγραμματισμένη ελάχιστη εφαρμογή περιλαμβάνει χάρτη πλήρους έργου, \(D_{\mathrm{fam}}\) και \(D_{\mathrm{within}}\). Η εκτεταμένη εφαρμογή προσθέτει \(D_{\mathrm{seq}}\) και \(D_{\mathrm{func}}\) μόνο όταν επαρκούν τα δεδομένα διαδοχής, λειτουργικής τοποθέτησης και σύγκρισης ανεξάρτητων έργων.

\subsection*{Η περίπτωση του Νικολάου Πρωτοψάλτη Σμύρνης}
Ο ιστορικός έλεγχος εξετάζει αν και πού διαφοροποιείται το έργο του Νικολάου ως προς το δηλωμένο περιβάλλον αναφοράς. Τα έργα του δεν χρησιμοποιούνται για εκπαίδευση του συστήματος στο αντίστοιχο τελικό ιστορικό ερώτημα. Η διαφορά μπορεί να αφορά επιλογή οικογενειών, διαδοχή, λειτουργική τοποθέτηση ή εσωτερική ανάπτυξη. Καθεμία ερμηνεύεται χωριστά και αναφέρεται σε συγκεκριμένα μουσικά χωρία.

\subsection*{Όσα δεν προκύπτουν από τα προφίλ}
Η αριθμητική απόκλιση δεν αποτελεί απόδειξη ιστορικής καινοτομίας, αισθητικής ανωτερότητας, λάθους στη σύνθεση ή απομάκρυνσης από μια αμετάβλητη παράδοση. Η μουσικολογική απόδοση τέτοιων χαρακτηρισμών απαιτεί ανεξάρτητη κριτική εξέταση των πηγών και σαφώς ορισμένα ερωτήματα.


% 15
\section{Στατιστική πειθαρχία και έλεγχος αξιοπιστίας}
\subsection*{Χωρισμός των δεδομένων ανά ερευνητικό ρόλο}
Τα δεδομένα εκπαίδευσης χρησιμεύουν για την εκμάθηση των παραμέτρων. Τα δεδομένα ανάπτυξης χρησιμεύουν για την επιλογή αρχιτεκτονικής, αναπαράστασης και υπερπαραμέτρων. Τα δεδομένα τελικής αξιολόγησης παραμένουν χωριστά ώσπου να ολοκληρωθούν οι προηγούμενες αποφάσεις. Το ίδιο έργο ή εξαρτημένες πηγές του δεν πρέπει να εμφανίζονται σε διαφορετικούς διαχωρισμούς.

\subsection*{Η ανάγκη για αντιστοιχισμένες συγκρίσεις}
Στη M1 διατηρείται κοινός στόχος αξιολόγησης. Στη M2 οι B2 και EISE συγκρίνονται με ίδια αρχιτεκτονική και αντίστοιχους πόρους, αλλά με διαφορετική προέλευση της σχεσιακής εποπτείας. Στη M3b ελέγχονται οι συμβολικές και ακουστικές αποστάσεις. Στη M3c τα δύο πληροφοριακά μεγέθη εκτιμώνται για την ίδια μουσική συνέχεια. Οι αναλύσεις αυτές επιδιώκουν να περιορίσουν εύλογες εναλλακτικές εξηγήσεις.

\subsection*{Αβεβαιότητα και πλήθος ανεξάρτητων παρατηρήσεων}
Η εκτίμηση της αβεβαιότητας πρέπει να λαμβάνει υπόψη ότι εμφανίσεις του ίδιου έργου και συγγενείς πηγές δεν είναι ανεξάρτητες. Το ΠΕ0 προσδιορίζει το πλήθος ανεξάρτητων έργων, πηγών, οικογενειών, ιστορικών σχέσεων και εκτελεστών και υποστηρίζει τις προσομοιώσεις απαιτούμενου δείγματος. Όπου η κάλυψη είναι ανεπαρκής, η αντίστοιχη ανάλυση παραμένει περιγραφική ή διερευνητική.

\subsection*{Κριτήρια ενεργοποίησης}
Οι προκαθορισμένες συνθήκες GREEN, AMBER και RED περιγράφουν αντίστοιχα αν μια μέτρηση επαρκεί στην προτιμώμενη κλίμακα, μόνο σε αδρότερη μουσική κλίμακα ή δεν μπορεί να στηρίξει αξιόπιστο κύριο συμπέρασμα. Οι απαιτήσεις δεν αναθεωρούνται εκ των υστέρων για να βελτιωθεί ένα αποτέλεσμα.

\subsection*{Όταν δεν υπάρχει επιλέξιμη αναπαράσταση}
Αν \(\widehat{\mathcal A}_{\mathrm{dec}}=\varnothing\), δεν δηλώνεται τεχνητά μια επιτυχής \(R^\star\). Επιτρέπονται περιορισμένοι διαγνωστικοί έλεγχοι για να εντοπιστεί το πρόβλημα, αλλά δεν υποκαθιστούν την προβλεπόμενη επιβεβαιωτική αξιολόγηση της M2.


% 16
\section{Οργάνωση της διατριβής και χρονική ακολουθία}
\subsection*{Βασικός κορμός}
Ο υπολογιστικός και μουσικολογικός πυρήνας συγκροτείται από την επιλογή αναπαράστασης (M1), την ελεγχόμενη σύγκριση B2--EISE (M2), την ανεξάρτητη επικύρωση του λεξιλογίου (M3a) και το αναγκαίο εκτελεστικό κριτήριο της M4. Ο ιστορικός χαρακτηρισμός M5 έχει συγκεκριμένο ελάχιστο παραδοτέο. Οι ανθρώπινες δοκιμασίες M3c και M3b διαθέτουν ιδιαίτερες προϋποθέσεις επάρκειας, ενώ οι νευρικές μετρήσεις αποτελούν επέκταση μόνο υπό επιβεβαιωμένη δυνατότητα υλοποίησης.

\subsection*{Ενδεικτική πορεία σαράντα οκτώ μηνών}
Στους πρώτους μήνες προηγούνται η απογραφή του υλικού, η εξασφάλιση αδειών, η συμβολική μετατροπή των πηγών, οι πρώτες οικογενειακές κάρτες και ο έλεγχος επισημείωσης. Η M1 οδηγεί στην επιλογή της \(R^\star\). Η M3a μπορεί να εξελίσσεται παράλληλα με άλλες εργασίες, καθώς το υλικό της παραμένει ανεξάρτητο. Η M2 εξετάζει την ιστορική εποπτεία με κατάλληλο εξωτερικό εκτελεστικό κριτήριο. Η M3c προηγείται της M3b στη συλλογή κρίσεων από τους ίδιους συμμετέχοντες. Η M4 επεκτείνεται όπου η ευθυγράμμιση επιτρέπει αξιόπιστες μετρήσεις και η M5 καταλήγει στην ιστορική εφαρμογή.

\subsection*{Οριστικοποιημένες αποφάσεις και τεχνικά ανοικτά σημεία}
Στην πρόταση είναι καθορισμένα τα βασικά ερευνητικά ερωτήματα, η διασταυρωμένη διάταξη \(U\times F\), η λογική επιλογής της \(R^\star\), η διάκριση B2--EISE, η ανάγκη ανεξάρτητης επικύρωσης και η ελάχιστη εφαρμογή της M5. Παραμένουν ανοικτά το πραγματικό πλήθος ανεξάρτητων πηγών, η αποτελεσματική συμβολική μεταγραφή, η λεπτομέρεια ηχητικής ευθυγράμμισης, η επάρκεια ακροαματικών ερεθισμάτων και οι τελικές παράμετροι των μοντέλων.

\subsection*{Όρια της παρούσας εκδοχής}
Η v8.71a διατηρεί τις M3b και M3c ως επιβεβαιωτικές ανθρώπινες δοκιμασίες υπό συγκεκριμένα κριτήρια ενεργοποίησης. Η πιθανή μεταφορά της M3b σε περισσότερο περιορισμένη βαθμίδα αποτελεί υπό συζήτηση αναθεώρηση και δεν περιγράφεται ως ήδη υιοθετημένη επιλογή.


% PAGE 18
\section{Κριτήρια επιστημονικής επάρκειας}

\subsection*{Η επίδοση ενός μοντέλου δεν αρκεί από μόνη της}
Η αναγνώριση μουσικών θέσεων αξιολογείται σε περισσότερα από ένα επίπεδα. Η ακριβής οριοθέτηση δεν εγγυάται ορθή απόδοση οικογένειας. Η σωστή οικογενειακή απόδοση δεν εγγυάται ότι δύο πραγματώσεις συνδέονται με ιστορική εξήγηση. Τέλος, μια υπολογιστικά αξιόπιστη κατηγορία δεν είναι υποχρεωτικό να επηρεάζει την ανθρώπινη μουσική προσδοκία. Οι διαφορετικές αυτές αξιώσεις χρειάζονται αντίστοιχες ανεξάρτητες μετρήσεις.

\subsection*{Η αρχή των εξωτερικών τεκμηρίων}
Η M1 μετρά τα όρια και τις οικογενειακές αποδόσεις σε υλικό με γνωστή προέλευση. Η M2 ελέγχει αν η ιστορική εποπτεία προσφέρει πρόσθετη πληροφορία σε σχέση με μια δίκαια αντιστοιχισμένη μη ιστορική συνθήκη. Η M3a εξετάζει την αξιοπιστία ανεξάρτητων μουσικολογικών κρίσεων. Οι M3c και M3b εξετάζουν διαφορετικές διαστάσεις της σημερινής ακρόασης, ενώ η M4 προσθέτει μετρήσεις πραγματικών εκτελέσεων. Η M5 αξιοποιεί όσα έχουν επιβεβαιωθεί για τον χαρακτηρισμό έργων.

\subsection*{Τρία επίπεδα τεκμηρίωσης}
\begin{enumerate}
\item \textbf{Ιστορικό τεκμήριο}: χειρόγραφο, έντυπη πηγή, καταγεγραμμένη εξήγηση, γνωστό έργο και επαληθευμένη σχέση πηγών.
\item \textbf{Μουσικολογική κρίση}: οριοθέτηση, οικογενειακή ένταξη, τροπική λειτουργία, χαρακτηρισμός παραλλαγής και τεκμηριωμένη αβεβαιότητα.
\item \textbf{Υπολογιστικό αποτέλεσμα}: πιθανότητα, διανυσματική αναπαράσταση, δείκτης διαφοροποίησης ή προβλεπτική επίδοση.
\end{enumerate}
Κανένα επίπεδο δεν πρέπει να υποκαθιστά τα άλλα. Ιδίως, μια υψηλή ομοιότητα που υπολογίζει το μοντέλο δεν μετατρέπεται σε ιστορικό τεκμήριο.

\subsection*{Αρνητικά αποτελέσματα και περιορισμός ισχυρισμών}
Η αποτυχία μιας επιβεβαιωτικής υπόθεσης αναφέρεται ως τέτοια. Αν η M3c δεν δείξει πρόσθετη συμβολή της οικογενειακής πληροφορίας, περιορίζεται η συγκεκριμένη γνωσιακή αξίωση και όχι αυτομάτως η μουσικολογική υπόσταση της θέσης. Αν η EISE δεν υπερέχει της B2, η ιστορική εποπτεία δεν τεκμηριώνεται ότι προσθέτει αξία υπό το συγκεκριμένο κριτήριο. Αν οι ειδικοί της M3a διαφωνούν έντονα, απαιτείται επανεξέταση της αξιοπιστίας του λεξιλογίου. Δεν διατυπώνεται γενικό συμπέρασμα από μεμονωμένη υποομάδα.

% PAGE 19
\section{Δυσκολίες των δεδομένων και περιορισμοί της μεθόδου}

\subsection*{Μουσικολογικά αβέβαιες εμφανίσεις}
Ορισμένα χωρία έχουν εναλλακτικά αποδεκτά όρια, συνδετικό υλικό ή περισσότερες από μία πιθανές οικογενειακές αποδόσεις. Η καταγραφή αυτή αποτελεί πληροφορία για το ίδιο το μουσικό αντικείμενο και όχι απλή ατέλεια της επισημείωσης. Η εκπαίδευση με μερικές ετικέτες επιτρέπει πολλαπλές συμβατές αναλύσεις, ενώ η τελική αξιολόγηση διαχωρίζει τα σφάλματα εντοπισμού από τις ουσιαστικές αμφισημίες.

\subsection*{Ανεπαρκή ιστορικά ζεύγη}
Η ύπαρξη πολλών εμφανίσεων μιας οικογένειας δεν συνεπάγεται ότι υπάρχουν εξίσου πολλά ανεξάρτητα ζεύγη ιστορικής εξήγησης. Κάθε τέτοια σχέση απαιτεί πηγολογικό έλεγχο. Όταν τα ιστορικά ζεύγη είναι λίγα, η κοινή προεκπαίδευση σε ευρύτερο συμβολικό σώμα μπορεί να βοηθήσει τις βασικές αναπαραστάσεις, αλλά δεν αυξάνει τον αριθμό των ανεξάρτητων ιστορικών τεκμηρίων. Ελεγχόμενα τεχνητά παραδείγματα δεν παρουσιάζονται ως ιστορικές εξηγήσεις.

\subsection*{Κίνδυνος εξάρτησης των δεδομένων}
Αν συγγενείς εκδοχές του ίδιου έργου βρεθούν και στην εκπαίδευση και στην τελική αξιολόγηση, η μετρημένη επίδοση μπορεί να είναι υπερβολικά αισιόδοξη. Οι διαχωρισμοί γίνονται επομένως σε επίπεδο έργου και ομάδας συγγενών πηγών, με ειδικούς αποκλεισμούς όπου το ερώτημα απαιτεί διαχωρισμό μελοποιών. Η παραγωγή τεχνητών παραλλαγών ακολουθεί τους διαχωρισμούς του πρωτοτύπου και δεν τους παρακάμπτει.

\subsection*{Η δυσκολία της ηχητικής ευθυγράμμισης}
Στην πραγματική εκτέλεση μία γραπτή συλλαβή μπορεί να περιέχει εκτενή μελωδική ανάπτυξη, συντομεύσεις, πρόσθετες κινήσεις ή καταληκτικές επιμηκύνσεις. Η χρονική ευθυγράμμιση από μόνη της δεν διαχωρίζει όλες αυτές τις αιτίες. Γι' αυτό προβλέπονται χειροκίνητα ελεγμένα σημεία αναφοράς, μέτρηση σφάλματος και επιλογή μεγαλύτερης μουσικής μονάδας όταν η λεπτομερής αντιστοίχιση δεν είναι αξιόπιστη.

\subsection*{Επαρκές υλικό για τις ανθρώπινες δοκιμασίες}
Οι ακροαματικές δοκιμασίες απαιτούν μουσικά φυσικές συνέχειες που διαφοροποιούν τα υπό εξέταση πληροφοριακά μεγέθη και τριάδες με συγκρίσιμες επιφανειακές αποστάσεις. Αν η αναγκαία κάλυψη λείπει, δεν δημιουργούνται αυθαίρετα παραδείγματα για να διατηρηθεί ένας επιβεβαιωτικός ισχυρισμός. Προηγείται περιορισμένος έλεγχος επάρκειας του υλικού.

% PAGE 20
\section{Μουσική θέση, θεωρία και ιστορική ερμηνεία}

\subsection*{Η ιστορική θεωρία ως αφετηρία και όχι ως έτοιμος αλγόριθμος}
Η αναφορά του Μανουήλ Χρυσάφη στη συνένωση μουσικών σημείων και η αναλογία του Γαβριήλ ιερομονάχου με τις λέξεις της γραμματικής υποστηρίζουν την ανάγκη μιας ενδιάμεσης μουσικής μονάδας. Δεν ορίζουν, όμως, αυτομάτως τα σημερινά όρια οικογενειών ή την κατάλληλη υπολογιστική αναπαράσταση. Οι αναφορές στον Χρύσανθο για τον ρόλο της βαθμίδας και του μουσικού περιβάλλοντος προσφέρουν ιστορικό κίνητρο για εξέταση του συμφραζομένου, όχι απόδειξη ότι οι θέσεις ακολουθούν το πιθανοτικό υπόδειγμα της M3c.

\subsection*{Η εξήγηση ως τεκμηριωμένη σχέση}
Η ιστορική εξήγηση αφορά συγκεκριμένες πηγές και αναλυτικές αποδόσεις. Δύο πραγματώσεις μπορεί να είναι συγγενείς χωρίς να διαθέτουν γνωστή εξηγητική σχέση. Μπορεί επίσης μια τεκμηριωμένη εξηγητική απόδοση να είναι επιφανειακά πολύ διαφορετική. Η ιστορική προέλευση κάθε σχέσης πρέπει να διατηρείται χωριστά από τις υπολογιζόμενες μελωδικές αποστάσεις.

\subsection*{Εκτέλεση και επιτέλεση}
Η εκτέλεση αφορά την ακουστή πραγμάτωση του μέλους. Η επιτέλεση περιλαμβάνει το ευρύτερο τελετουργικό και κοινωνικό περιβάλλον: τη λειτουργική περίσταση, τον χώρο, τους συμμετέχοντες και τις πρακτικές μετάδοσης. Η M4 μετρά συγκεκριμένες πλευρές της εκτέλεσης. Η μεταβλητή \(R\) της M3c είναι περιορισμένη πειραματική ένδειξη λατρευτικού συμφραζομένου· δεν αναπαριστά ολόκληρη τη λατρευτική επιτέλεση.

\subsection*{Σύγχρονη ακρόαση και ιστορική μεταβλητότητα}
Οι σημερινές ακροαματικές κρίσεις δεν ανασυνθέτουν την εμπειρία ακροατών προηγούμενων αιώνων. Η ιστορική ακρόαση προσεγγίζεται μόνο έμμεσα μέσα από πηγές, περιγραφές πρακτικών και προσεκτική ιστορική ερμηνεία. Η σύγχρονη ακρόαση εξετάζεται πειραματικά ως διαφορετικό είδος τεκμηρίου. Η συνδυαστική ερμηνεία τους απαιτεί ρητή διατήρηση αυτής της διάκρισης.

% PAGE 21
\section{Σύνοψη της ερευνητικής αρχιτεκτονικής}

\subsection*{Από την παρτιτούρα στον χαρακτηρισμό του έργου}
Η αρχική πηγή μετατρέπεται σε συμβολικά μελωδικά γεγονότα με διατηρημένη προέλευση. Η M1 επιλέγει την κατάλληλη αναπαράσταση για αναγνώριση ορίων και οικογενειών. Η M2 χρησιμοποιεί την επιλεγμένη αναπαράσταση, εντοπίζει υποψήφιες εμφανίσεις και εξετάζει τη συμβολή τεκμηριωμένων ιστορικών εξηγητικών σχέσεων. Η M3a παρέχει ανεξάρτητη μουσικολογική αξιολόγηση και οι M3c/M3b εξετάζουν, υπό τους αντίστοιχους όρους εφικτότητας, τη μουσική προσδοκία και την ακροαματική κατηγοριοποίηση. Η M4 συνδέει τη γραπτή πραγματοποίηση με τη σύγχρονη εκτέλεση. Η M5 χαρακτηρίζει ολόκληρα έργα ως προς ορισμένο σώμα αναφοράς.

\subsection*{Οι βασικοί συμβολισμοί}
\begin{tabularx}{\linewidth}{>{\bfseries}p{2.3cm}X>{\bfseries}p{2.5cm}X}
\(R^\star\) & Επιλεγμένη αναπαράσταση & \(z_{\mathrm{id}}\) & Οικογενειακή ταυτότητα\\
\(\mathcal R^\star\) & Αβεβαιότητα επιλογής & \(h_t\) & Τοπικές καταστάσεις\\
\(g_{\min}\) & Ελάχιστη οικογενειακή επάρκεια & \(\Delta_{\mathrm{hist}}\) & EISE μείον B2\\
\(M(R)\) & Επίδοση ορίων & \(I_F\) & Απρόσμενο οικογένειας\\
\(G(R)\) & Επίδοση αναγνώρισης οικογενειών & \(D_{\mathrm{within}}\) & Διαφορά πραγματώσεων\\
\end{tabularx}

\subsection*{Ο εννοιολογικός πυρήνας}
Η μουσική θέση αντιμετωπίζεται ως μουσικολογικά ορισμένη και υπολογιστικά ελέγξιμη μονάδα μεταβλητής έκτασης. Η ύπαρξη οικογενειακής συγγένειας δεν συνεπάγεται ιστορική εξηγητική σχέση. Η επίδοση αξιολογείται σε πραγματικά χωρισμένα δεδομένα. Η σύγχρονη ανθρώπινη ακρόαση εξετάζεται ανεξάρτητα από τις ιστορικές πηγές. Η διαφοροποίηση ενός έργου περιγράφεται ως προς ρητό περιβάλλον αναφοράς, χωρίς να μετατρέπεται σε αισθητική βαθμολογία.

\begin{keybox}
\textbf{Κεντρικός σκοπός.} Να συγκροτηθεί ένα συνεκτικό πλαίσιο που ενώνει την αναγνώριση μουσικών θέσεων, τη μελέτη της ιστορικής παραλλαγής και τον πολυδιάστατο χαρακτηρισμό πλήρων έργων, με διακριτούς ελέγχους μουσικολογικής, υπολογιστικής και εκτελεστικής εγκυρότητας.
\end{keybox}

% 21
\section{Από τη μουσική πηγή στο αποτέλεσμα του συστήματος}
\subsection*{Η διαδρομή μιας εμφάνισης}
\begin{enumerate}
\item Εξαγωγή της μουσικής πληροφορίας από πηγή της Νέας Μεθόδου με διατήρηση έργου, σελίδας, χωρίου και προέλευσης ανά συμβολικό γεγονός.
\item Χρήση της αναπαράστασης \(R^\star\), όπως προέκυψε από τη M1, για τη συγκρότηση της εισόδου.
\item Υπολογισμός τοπικής αριθμητικής κατάστασης \(h_t\) ανά γεγονός από τον κωδικοποιητή και σχηματισμός επιλέξιμων συνεχόμενων τμημάτων.
\item Συμπύκνωση των καταστάσεων κάθε τμήματος σε \(z_{\mathrm{id}}\) για την οικογενειακή αναγνώριση, με χωριστή διατήρηση πληροφορίας της εσωτερικής πραγμάτωσης.
\item Κοινή επιλογή ορίων και οικογενειών, με δυνατότητα ανεπίλυτης απόδοσης, αποχής ή αγνώστου οικογένειας.
\item Παραγωγή χάρτη εμφανίσεων για το έργο: όρια, οικογένεια, αβεβαιότητα, μη καλυμμένες περιοχές και συγκρίσιμες πραγματώσεις από τεκμηριωμένες πηγές.
\end{enumerate}

\subsection*{Ο ρόλος της ιστορικής εποπτείας}
Η B2 και η EISE χρησιμοποιούν την ίδια αρχιτεκτονική και το ίδιο λεξιλόγιο ανάπτυξης. Η κύρια διαφορά είναι ότι τα θετικά ζεύγη της B2 προέρχονται από κοινή οικογενειακή ένταξη, ενώ τα αντίστοιχα ζεύγη της EISE απαιτούν ανεξάρτητα επαληθευμένες ιστορικές εξηγητικές σχέσεις. Η επαλήθευση αυτή προηγείται της εκπαίδευσης και δεν παράγεται από το μοντέλο.

\subsection*{Προεκπαίδευση και ελεγχόμενα τεχνητά δεδομένα}
Μπορεί να αξιοποιηθεί κοινή προεκπαίδευση σε ευρύτερο συμβολικό σώμα, με εργασίες όπως η συμπλήρωση αποκρυμμένων μελωδικών γεγονότων. Ελεγχόμενες τεχνητές παραλλαγές μπορούν να βοηθήσουν τη μάθηση όταν επιτρέπεται μουσικολογικά, αλλά δεν αποτελούν ιστορικά τεκμήρια. Οι παράγωγες παραλλαγές ενός πρωτοτύπου ακολουθούν τον ίδιο διαχωρισμό εκπαίδευσης ή αξιολόγησης, ώστε να αποφεύγεται η διαρροή πληροφορίας.


% 22
\section{Βασικές διακρίσεις που πρέπει να διατηρούνται}
\begin{enumerate}
\item \textbf{Νεύμα και γεγονός.} Ένας σημειογραφικός χαρακτήρας δεν αντιστοιχεί υποχρεωτικά σε ένα ατομικό μελωδικό γεγονός.
\item \textbf{Τμήμα και θέση.} Δεν είναι κάθε συνεχόμενο τμήμα μουσική θέση. Η αναγνώριση απαιτεί μουσικολογικά και υπολογιστικά κριτήρια.
\item \textbf{Ομοιότητα και οικογένεια.} Η μελωδική ομοιότητα αποτελεί ένδειξη, όχι από μόνη της απόδειξη οικογενειακής ένταξης.
\item \textbf{Οικογένεια και ιστορική εξήγηση.} Η τεκμηριωμένη εξηγητική σχέση είναι ιδιαίτερο γεγονός της ιστορίας των πηγών.
\item \textbf{Άγνωστη οικογένεια και νέος τύπος.} Η αποχή του μοντέλου δεν αποδεικνύει ιστορική ανακάλυψη.
\item \textbf{Τυφλή επισημείωση.} Οι ανεξάρτητοι ειδικοί διαθέτουν το απαραίτητο μουσικό συμφραζόμενο, όχι τις προβλέψεις του μοντέλου ή τις μεταξύ τους κρίσεις.
\item \textbf{Κατηγοριοποίηση και ομοιότητα.} Η M3b εξετάζει αν η οικογενειακή συγγένεια προσθέτει πληροφορία πέρα από τις μελωδικές και ακουστικές αποστάσεις.
\item \textbf{Προσδοκία οικογένειας και γεγονότος.} Το \(I_F\) και το \(I_{\mathrm{event}}\) προέρχονται από διαφορετικές κλίμακες πρόβλεψης και χρειάζονται χωριστή αξιολόγηση.
\item \textbf{Αλγόριθμος και πρόσθετη πληροφορία.} Η χρήση προηγούμενου συμφραζομένου από το PPM δεν αποδεικνύει από μόνη της προβλεπτικό πλεονέκτημα. Απαιτείται σύγκριση σε ανεξάρτητα έργα.
\item \textbf{Πειραματική ένδειξη και λατρευτική επιτέλεση.} Η μεταβλητή \(R\) της M3c δεν αναπαριστά την πλήρη εμπειρία μιας ακολουθίας.
\item \textbf{Εκτελεστική ομοιότητα και ιστορική μετάδοση.} Η σύγχρονη πραγμάτωση μπορεί να μοιάζει με γνωστή εξήγηση χωρίς να τεκμηριώνεται άμεση ιστορική συνέχεια.
\item \textbf{Χαρακτηρισμός έργου και αισθητική αξιολόγηση.} Οι δείκτες της M5 περιγράφουν διαφοροποίηση από ορισμένο σώμα αναφοράς, όχι μουσική αξία ή ποιότητα.
\end{enumerate}
\subsection*{Συνολική αρχή}
Ο μουσικολογικός ισχυρισμός πρέπει να παραμένει ανάλογος προς το τεκμήριο: οι πηγές θεμελιώνουν ιστορικές σχέσεις, οι ανεξάρτητοι ειδικοί ελέγχουν την αναγνώριση, οι υπολογιστικές δοκιμασίες μετρούν προβλεπτική και αναπαραστατική ικανότητα, ενώ οι σύγχρονες ακροάσεις και εκτελέσεις εξετάζουν συγκεκριμένες πλευρές της μουσικής εμπειρίας.


\section{Συνοπτικό γλωσσάρι}
\small
\begin{longtable}{>{\bfseries}p{4.0cm}p{10.7cm}}
\toprule
Έννοια & Επεξήγηση \\
\midrule
\endfirsthead
\toprule Έννοια & Επεξήγηση \\ \midrule
\endhead
Θέση & Αναγνωρίσιμος μελικός σχηματισμός που μπορεί να αποκτά διαφορετική έκταση και πραγμάτωση. \\
Οικογένεια θέσεων & Ομάδα συγγενών πραγματώσεων με ορισμένα διαγνωστικά μουσικά γνωρίσματα. \\
Εμφάνιση & Συγκεκριμένο χωρίο ενός έργου στο οποίο αναγνωρίζεται πιθανή θέση. \\
Πραγμάτωση & Η συγκεκριμένη μορφή με την οποία παρουσιάζεται μια θέση σε ένα χωρίο. \\
Ιστορική εξήγηση & Τεκμηριωμένη αναλυτική απόδοση παλαιότερης μουσικής γραφής σε νεότερη μορφή. \\
Συμβολική αναπαράσταση & Κωδικοποίηση της μουσικής πληροφορίας σε μονάδες που μπορεί να επεξεργαστεί υπολογιστικά ένα μοντέλο. \\
Ατομικό μελωδικό γεγονός & Ελάχιστη μονάδα μελωδικής πληροφορίας του επιλεγμένου σχήματος, όχι υποχρεωτικά ένα νεύμα. \\
Τμηματοποίηση & Εντοπισμός των ορίων μουσικών ενοτήτων μέσα σε ακολουθία γεγονότων. \\
Κωδικοποιητής & Μέρος μοντέλου που μετατρέπει τα συμβολικά γεγονότα σε αριθμητικές αναπαραστάσεις. \\
Πρότυπο οικογένειας & Χαρακτηριστική αριθμητική αναπαράσταση που χρησιμοποιείται στη σύγκριση συγγενών πραγματώσεων. \\
Μερική επισημείωση & Κατάσταση στην οποία ορισμένες εμφανίσεις ή οικογενειακές ετικέτες δεν έχουν προσδιοριστεί με ασφάλεια. \\
Αποχή & Απόφαση του μοντέλου να μην αποδώσει ανασφαλή οικογενειακή ταυτότητα. \\
Εκπαίδευση & Προσαρμογή των παραμέτρων του μοντέλου στα διαθέσιμα παραδείγματα. \\
Ανάπτυξη & Χρήση χωριστών δεδομένων για επιλογή μοντέλου και παραμέτρων. \\
Τελική αξιολόγηση & Έλεγχος σε δεδομένα που δεν χρησιμοποιήθηκαν για την επιλογή του μοντέλου. \\
Ανάκληση & Ποσοστό των πραγματικών περιπτώσεων που εντοπίζονται σωστά. \\
Ακρίβεια & Ποσοστό των θετικών προβλέψεων που είναι σωστές. \\
Δείκτης F1 & Συνοπτικό μέτρο που συνδυάζει την ακρίβεια και την ανάκληση. \\
EISE & Σύστημα μάθησης που χρησιμοποιεί μουσικές οικογένειες και επαληθευμένες ιστορικές εξηγητικές σχέσεις. \\
Ημιμαρκοβιανό CRF & Πιθανοτικό μοντέλο που βαθμολογεί συνεχόμενα τμήματα μεταβλητού μήκους και τις ετικέτες τους. \\
PPM & Μέθοδος πρόβλεψης από προηγούμενο συμφραζόμενο μεταβλητού μήκους, με επιστροφή σε μικρότερο όταν τα δεδομένα είναι αραιά. \\
Πληροφοριακό περιεχόμενο & Αριθμητικό μέγεθος που αυξάνεται όταν μια μουσική επιλογή είναι λιγότερο πιθανή κατά το μοντέλο. \\
\(I_F\) & Πληροφοριακό περιεχόμενο της επόμενης οικογένειας μουσικών θέσεων. \\
\(I_{\mathrm{event}}\) & Πληροφοριακό περιεχόμενο των επιμέρους μελωδικών γεγονότων της συνέχειας. \\
Ευθυγράμμιση & Χρονική αντιστοίχιση των γραπτών μουσικών γεγονότων με μια ηχογράφηση. \\
\(A(z)\) & Προφίλ ιστορικής μεταβλητότητας σε εσωτερική περιοχή μιας θέσης. \\
\(\Gamma(z)\) & Προφίλ μεταβλητότητας μεταξύ σύγχρονων εκτελεστών στην αντίστοιχη περιοχή. \\
\(\Omega\) & Δηλωμένο σώμα πηγών και κανόνων που χρησιμεύει ως περιβάλλον αναφοράς για τον χαρακτηρισμό έργου. \\
\(R^\star\) & Η αναπαράσταση που επιλέγει η M1 για την κύρια σύγκριση της M2. \\
\(M(R)\) & Επίδοση εντοπισμού ορίων μιας υποψήφιας αναπαράστασης. \\
\(G(R)\) & Επίδοση αναγνώρισης οικογενειών μιας υποψήφιας αναπαράστασης. \\
\(g_{\min}\) & Προκαθορισμένο κατώτατο αποδεκτό επίπεδο οικογενειακής επάρκειας. \\
B0 & Άμεσος συμβολικός χαρακτηρισμός ιστορικής παραλλαγής, χωρίς μάθηση αναπαράστασης. \\
B1 & Κοινή αρχιτεκτονική χωρίς πρόσθετη εποπτεία από ζεύγη πραγματώσεων. \\
B2 & Εκπαίδευση με ζεύγη κοινής οικογένειας, χωρίς τεκμηριωμένη εξηγητική σχέση. \\
\(\Delta_{\mathrm{hist}}\) & Διαφορά συμφωνίας EISE και B2 με το ίδιο ανεξάρτητο εκτελεστικό κριτήριο. \\
\(E_{\mathrm{Byz}}\) & Διαβαθμισμένος δείκτης εξοικείωσης του ακροατή με τη βυζαντινή μουσική. \\
\(R\) στη M3c & Πειραματική, μη μουσική ένδειξη λατρευτικού συμφραζομένου. \\
\(D_{\mathrm{fam}}\) & Διαφοροποίηση στη χρήση και στην κατανομή οικογενειών μέσα σε έργο. \\
\(D_{\mathrm{seq}}\) & Διαφοροποίηση στη διαδοχή των οικογενειών. \\
\(D_{\mathrm{func}}\) & Διαφοροποίηση στη δομική ή λειτουργική τοποθέτηση οικογενειών. \\
\(D_{\mathrm{within}}\) & Διαφοροποίηση στον τρόπο πραγμάτωσης της ίδιας οικογένειας. \\
\bottomrule
\end{longtable}
\normalsize

\end{document}

